%%%%%%%%%%%%%%%%
%%% exod 10 %%%
%%%%%%%%%%%%%%%%

\Chapter{10}

\Verse{1}
{
	\hRJE{וַיֹּאמֶר יְהֹוָה אֶל מֹשֶׁה בֹּא אֶל פַּרְעֹה כִּי אֲנִי הִכְבַּדְתִּי אֶת לִבּוֹ וְאֶת לֵב עֲבָדָיו לְמַעַן שִׁתִי אֹתֹתַי אֵלֶּה בְּקִרְבּוֹ׃}
}
{
	\eRJE{
		And YHWH said to Moses, “Come to Pharaoh, because I have
		made his heart and his servants’ heart heavy for the
		purpose of my setting these signs of mine among them
	}
}
{
%	\footnote{
%		I have made his heart heavy. Earlier the text says that Pharaoh made his
%		heart heavy. Now it says that God claims to have done it. This verse may
%		refer only to the immediate condition; that is, that Pharaoh hardened
%		his own heart before, but now it is God doing it. Or it may suggest dual
%		causation; that is, even when the human (Pharaoh) decides himself, it is
%		still God who is causing this decision. It is in the theme of YHWH’s
%		becoming known that we may find the explanation of the contradiction
%		that has occurred even to little children: why does God harden Pharaoh’s
%		heart against releasing Israel and then plague Egypt for not letting
%		them go?! The fact is that the Torah does not present the plagues as
%		punishment of the Egyptians for enslaving the Israelites, but rather as
%		signs by which YHWH will become known. Thus it pictures God stating
%		unambiguously here: I have made his heart and his servants’ heart heavy
%		for the purpose of my setting these signs of mine among them and for the
%		purpose that you will tell in the ears of your son and your son’s son
%		about how I abused Egypt and about my signs that I set among them, and
%		you will know that I am YHWH. Even more dramatically, YHWH tells Moses
%		to inform the Pharaoh that “And in fact I established you for this
%		purpose, for the purpose of showing you my power—and in order to tell my
%		name in all the earth” (9:16). Dispersed through the account are
%		frequent reminders that this is the purpose of the hardening of
%		Pharaoh’s heart (7:3–5; 11:9–10; 14:4,17–18) and that this is the
%		purpose of the plagues (7:17; 8:6,18; 9:14,29). The Egyptians’ suffering
%		in the plagues may be understood to be justified, a consequence of their
%		cruelty; but, still, the purpose of the plagues is not punishment but
%		rather to make YHWH known.
%	}
}

\Verse{2}
{
	\hRJE{וּלְמַעַן תְּסַפֵּר בְּאזְנֵי בִנְךָ וּבֶן בִּנְךָ אֵת אֲשֶׁר הִתְעַלַּלְתִּי בְּמִצְרַיִם וְאֶת אֹתֹתַי אֲשֶׁר שַׂמְתִּי בָם וִידַעְתֶּם כִּי אֲנִי יְהֹוָה׃}
}
{
	\eRJE{
		and for the purpose that you will tell in the ears of your
		son and your son’s son about how I abused Egypt and about
		my signs that I set among them, and you will know that I
		am YHWH.”
	}
}
{
}

\Verse{3}
{
	\hRJE{וַיָּבֹא מֹשֶׁה וְאַהֲרֹן אֶל פַּרְעֹה וַיֹּאמְרוּ אֵלָיו כֹּה אָמַר יְהֹוָה אֱלֹהֵי הָעִבְרִים עַד מָתַי מֵאַנְתָּ לֵעָנֹת מִפָּנָי שַׁלַּח עַמִּי וְיַעַבְדֻנִי׃}
}
{
	\eRJE{
		And Moses and Aaron came to Pharaoh and said to him,
		“YHWH, God of the Hebrews, said this: ‘How long do you
		refuse to be humbled in front of me? Let my people go, so
		they may serve me.
	}
}
{
}

\Verse{4}
{
	\hRJE{כִּי אִם מָאֵן אַתָּה לְשַׁלֵּחַ אֶת עַמִּי הִנְנִי מֵבִיא מָחָר אַרְבֶּה בִּגְבֻלֶךָ׃}
}
{
	\eRJE{
		Because if you refuse to let my people go, here, I’m
		bringing a locust swarm in your border tomorrow,
	}
}
{
}

\Verse{5}
{
	\hRJE{וְכִסָּה אֶת עֵין הָאָרֶץ וְלֹא יוּכַל לִרְאֹת אֶת הָאָרֶץ וְאָכַל אֶת יֶתֶר הַפְּלֵטָה הַנִּשְׁאֶרֶת לָכֶם מִן הַבָּרָד וְאָכַל אֶת כּל הָעֵץ הַצֹּמֵחַ לָכֶם מִן הַשָּׂדֶה׃}
}
{
	\eRJE{
		and it will cover the eye of the land, and one won’t be
		able to see the land! And it will eat the remains of what
		has survived that is left to you from the hail, and it
		will eat every tree that is growing for you from the
		field.
	}
}
{
%	\footnote{
%		cover the eye of the land. Meaning: the locust swarm will be so thick as
%		to make the land invisible. Footnote 10:5. the remains of what has
%		survived that is left. The phrase is a triple redundancy: yeter happlth
%		hanniš ’eret. Its purpose may be to develop the pun on the name of
%		Jether, Moses’ father-in-law (Exod 4:18).
%	}
}

\Verse{6}
{
	\hRJE{וּמָלְאוּ בָתֶּיךָ וּבָתֵּי כל עֲבָדֶיךָ וּבָתֵּי כל מִצְרַיִם אֲשֶׁר לֹא רָאוּ אֲבֹתֶיךָ וַאֲבוֹת אֲבֹתֶיךָ מִיּוֹם הֱיוֹתָם עַל הָאֲדָמָה עַד הַיּוֹם הַזֶּה וַיִּפֶן וַיֵּצֵא מֵעִם פַּרְעֹה׃}
}
{
	\eRJE{
		And they’ll fill your houses and the houses of all your
		servants and the houses of all Egypt, which your fathers
		and your fathers’ fathers did not see, from the day they
		were on the land until this day’.” And he turned and went
		out from Pharaoh.
	}
}
{
}

\Verse{7}
{
	\hRJE{וַיֹּאמְרוּ עַבְדֵי פַרְעֹה אֵלָיו עַד מָתַי יִהְיֶה זֶה לָנוּ לְמוֹקֵשׁ שַׁלַּח אֶת הָאֲנָשִׁים וְיַעַבְדוּ אֶת יְהֹוָה אֱלֹהֵיהֶם הֲטֶרֶם תֵּדַע כִּי אָבְדָה מִצְרָיִם׃}
}
{
	\eRJE{
		And Pharaoh’s servants said to him, “How long will this be
		a trap for us? Let the people go, so they may serve YHWH,
		their God. Don’t you know yet that Egypt has perished?!”
	}
}
{
%	\footnote{
%		Egypt has perished. They speak as if the country is already destroyed.
%		Interpreters and translators have tried to make their statement more
%		logical by understanding it to mean “Egypt is lost” or “ruined.” But I
%		think that Pharaoh’s servants are pictured here as speaking with
%		deliberate hyperbole. The Israelites themselves will use this very same
%		hyperbole later: following the Korah rebellion, they say, “We have
%		expired. We have perished” (Num 17:27).
%	}
}

\Verse{8}
{
	\hRJE{וַיּוּשַׁב אֶת מֹשֶׁה וְאֶת אַהֲרֹן אֶל פַּרְעֹה וַיֹּאמֶר אֲלֵהֶם לְכוּ עִבְדוּ אֶת יְהֹוָה אֱלֹהֵיכֶם מִי וָמִי הַהֹלְכִים׃}
}
{
	\eRJE{
		And Moses and Aaron were brought back to Pharaoh. And he
		said to them, “Go. Serve YHWH, your God. Who are the ones
		who are going?”
	}
}
{
}

\Verse{9}
{
	\hRJE{וַיֹּאמֶר מֹשֶׁה בִּנְעָרֵינוּ וּבִזְקֵנֵינוּ נֵלֵךְ בְּבָנֵינוּ וּבִבְנוֹתֵנוּ בְּצֹאנֵנוּ וּבִבְקָרֵנוּ נֵלֵךְ כִּי חַג יְהֹוָה לָנוּ׃}
}
{
	\eRJE{
		And Moses said, “We’ll go with our young and with our old,
		we’ll go with our sons and with our daughters, with our
		sheep and with our oxen, because we have a festival of
		YHWH.”
	}
}
{
}

\Verse{10}
{
	\hRJE{וַיֹּאמֶר אֲלֵהֶם יְהִי כֵן יְהֹוָה עִמָּכֶם כַּאֲשֶׁר אֲשַׁלַּח אֶתְכֶם וְאֶת טַפְּכֶם רְאוּ כִּי רָעָה נֶגֶד פְּנֵיכֶם׃}
}
{
	\eRJE{
		And he said to them, “YHWH would be with you like that,
		when I would let you and your infants go, see, because bad
		is in front of your faces.
	}
}
{
%	\footnote{
%		YHWH would be with you like that. It is difficult to convey the
%		Pharaoh’s cynicism that is expressed in the Hebrew jussive. He says “Let
%		it be so” in the sense that there is no chance at all that it could be
%		so. Footnote 10:10. and your infants. Pharaoh is still negotiating,
%		trying to give the minimum. First he refused to let the people have
%		their religious festival. His second position was that they could have
%		it in the land rather than leave Egypt. Now he says that they may leave,
%		but their children must stay—thus guaranteeing that the people will
%		return.
%	}
}

\Verse{11}
{
	\hRJE{לֹא כֵן לְכוּ נָא הַגְּבָרִים וְעִבְדוּ אֶת יְהֹוָה כִּי אֹתָהּ אַתֶּם מְבַקְשִׁים וַיְגָרֶשׁ אֹתָם מֵאֵת פְּנֵי פַרְעֹה׃}
}
{
	\eRJE{
		It is not like that. Go—the men!—and serve YHWH, because
		that is what you’re asking.” And he drove them out from
		Pharaoh’s face.
	}
}
{
%	\footnote{
%		the men. The women must stay as well. Footnote 10:11. he drove them out
%		from Pharaoh’s face. Meaning: from Pharaoh’s presence.
%	}
}

\Verse{12}
{
	\hRJE{וַיֹּאמֶר יְהֹוָה אֶל מֹשֶׁה נְטֵה יָדְךָ עַל אֶרֶץ מִצְרַיִם בָּאַרְבֶּה וְיַעַל עַל אֶרֶץ מִצְרָיִם וְיֹאכַל אֶת כּל עֵשֶׂב הָאָרֶץ אֵת כּל אֲשֶׁר הִשְׁאִיר הַבָּרָד׃}
}
{
	\eRJE{
		And YHWH said to Moses, “Reach out your hand at the land
		of Egypt for the locust, and it will come up on the land
		of Egypt and eat all the land’s vegetation, everything
		that the hail has left.”
	}
}
{
}

\Verse{13}
{
	\hRJE{וַיֵּט מֹשֶׁה אֶת מַטֵּהוּ עַל אֶרֶץ מִצְרַיִם וַיהֹוָה נִהַג רוּחַ קָדִים בָּאָרֶץ כּל הַיּוֹם הַהוּא וְכל הַלָּיְלָה הַבֹּקֶר הָיָה וְרוּחַ הַקָּדִים נָשָׂא אֶת הָאַרְבֶּה׃}
}
{
	\eRJE{
		And Moses reached out his staff at the land of Egypt, and
		YHWH drove an east wind through the land all that day and
		all the night. It was the morning: and the east wind had
		carried the locust swarm.
	}
}
{
}

\Verse{14}
{
	\hRJE{וַיַּעַל הָאַרְבֶּה עַל כּל אֶרֶץ מִצְרַיִם וַיָּנַח בְּכֹל גְּבוּל מִצְרָיִם כָּבֵד מְאֹד לְפָנָיו לֹא הָיָה כֵן אַרְבֶּה כָּמֹהוּ וְאַחֲרָיו לֹא יִהְיֶה כֵּן׃}
}
{
	\eRJE{
		And the locust swarm came up over all the land of Egypt
		and lingered in all of Egypt’s border, very heavy: before
		it there was no such locust swarm like it, and after it
		there will not be such.
	}
}
{
}

\Verse{15}
{
	\hRJE{וַיְכַס אֶת עֵין כּל הָאָרֶץ וַתֶּחְשַׁךְ הָאָרֶץ וַיֹּאכַל אֶת כּל עֵשֶׂב הָאָרֶץ וְאֵת כּל פְּרִי הָעֵץ אֲשֶׁר הוֹתִיר הַבָּרָד וְלֹא נוֹתַר כּל יֶרֶק בָּעֵץ וּבְעֵשֶׂב הַשָּׂדֶה בְּכל אֶרֶץ מִצְרָיִם׃}
}
{
	\eRJE{
		And it covered the eye of all the land, and the land was
		dark; and it ate all the land’s vegetation and the fruit
		of every tree that the hail had left, and not any plant
		was left, in the tree and in the field’s vegetation, in
		all the land of Egypt.
	}
}
{
}

\Verse{16}
{
	\hRJE{וַיְמַהֵר פַּרְעֹה לִקְרֹא לְמֹשֶׁה וּלְאַהֲרֹן וַיֹּאמֶר חָטָאתִי לַיהֹוָה אֱלֹהֵיכֶם וְלָכֶם׃}
}
{
	\eRJE{
		And Pharaoh hurried to call Moses and Aaron, and he said,
		“I’ve sinned against YHWH, your God, and against you.
	}
}
{
}

\Verse{17}
{
	\hRJE{וְעַתָּה שָׂא נָא חַטָּאתִי אַךְ הַפַּעַם וְהַעְתִּירוּ לַיהֹוָה אֱלֹהֵיכֶם וְיָסֵר מֵעָלַי רַק אֶת הַמָּוֶת הַזֶּה׃}
}
{
	\eRJE{
		And now, bear my sin just this one time and pray to YHWH,
		your God, that He’ll turn just this death away from me.”
	}
}
{
}

\Verse{18}
{
	\hRJE{וַיֵּצֵא מֵעִם פַּרְעֹה וַיֶּעְתַּר אֶל יְהֹוָה׃}
}
{
	\eRJE{And he went out from Pharaoh and prayed to YHWH,}
}
{
}

\Verse{19}
{
	\hRJE{וַיַּהֲפֹךְ יְהֹוָה רוּחַ יָם חָזָק מְאֹד וַיִּשָּׂא אֶת הָאַרְבֶּה וַיִּתְקָעֵהוּ יָמָּה סּוּף לֹא נִשְׁאַר אַרְבֶּה אֶחָד בְּכֹל גְּבוּל מִצְרָיִם׃}
}
{
	\eRJE{
		and YHWH turned back a very strong west wind, and it
		picked up the locust swarm and blew it to the Red Sea. Not
		one locust was left in all of Egypt’s border.
	}
}
{
%	\footnote{
%		Red Sea. Not the so-called “Reed” Sea. See the comment on Exod Footnote
%		13:18.
%	}
}

\Verse{20}
{
	\hRJE{וַיְחַזֵּק יְהֹוָה אֶת לֵב פַּרְעֹה וְלֹא שִׁלַּח אֶת בְּנֵי יִשְׂרָאֵל׃}
}
{
	\eRJE{
		And YHWH strengthened Pharaoh’s heart, and he did not let
		the children of Israel go.
	}
}
{
}

\Verse{21}
{
	\hRJE{וַיֹּאמֶר יְהֹוָה אֶל מֹשֶׁה נְטֵה יָדְךָ עַל הַשָּׁמַיִם וִיהִי חֹשֶׁךְ עַל אֶרֶץ מִצְרָיִם וְיָמֵשׁ חֹשֶׁךְ׃}
}
{
	\eRJE{
		And YHWH said to Moses, “Reach out your hand at the skies,
		and let there be darkness on the land of Egypt, and one
		will feel darkness.”
	}
}
{
%	\footnote{
%		one will feel darkness. Like the other plagues, this one is explicitly
%		identified as different from what is found in nature at other times. The
%		hail and locusts are unlike anything that has ever happened in Egypt.
%		The darkness can be felt. There is no chance whatever that the
%		Egyptians—or any interpreter of this text—can understand these as being
%		chance occurrences of nature. Note also that the word for “felt,” Hebrew
%		ymš, plays on the name Moses, mšeh.
%	}
}

\Verse{22}
{
	\hRJE{וַיֵּט מֹשֶׁה אֶת יָדוֹ עַל הַשָּׁמָיִם וַיְהִי חֹשֶׁךְ אֲפֵלָה בְּכל אֶרֶץ מִצְרַיִם שְׁלֹשֶׁת יָמִים׃}
}
{
	\eRJE{
		And Moses reached out his hand at the skies, and there was
		dismal darkness in all the land of Egypt—three days.
	}
}
{
%	\footnote{
%		three days. This is further indication that no natural occurrence, such
%		as an eclipse, is intended here; no eclipse lasts three days.
%	}
}

\Verse{23}
{
	\hRJE{לֹא רָאוּ אִישׁ אֶת אָחִיו וְלֹא קָמוּ אִישׁ מִתַּחְתָּיו שְׁלֹשֶׁת יָמִים וּלְכל בְּנֵי יִשְׂרָאֵל הָיָה אוֹר בְּמוֹשְׁבֹתָם׃}
}
{
	\eRJE{
		Each man did not see his brother, and each man did not get
		up from under it—three days. And for all the children of
		Israel there was light in their homes.
	}
}
{
%	\footnote{
%		from under it. The Hebrew () can mean “from under it” or “from his
%		place.” Everyone else makes it “from his place” or some equivalent, but
%		I think that the meaning “from under it” is possible, even preferable.
%		It conveys a much more vivid picture:
%	}
}

\Verse{24}
{
	\hRJE{וַיִּקְרָא פַרְעֹה אֶל מֹשֶׁה וַיֹּאמֶר לְכוּ עִבְדוּ אֶת יְהֹוָה רַק צֹאנְכֶם וּבְקַרְכֶם יֻצָּג גַּם טַפְּכֶם יֵלֵךְ עִמָּכֶם׃}
}
{
	\eRJE{
		And Pharaoh called Moses, and he said, “Go. Serve YHWH.
		Only your sheep and your oxen will stay put. Your infants
		will go with you as well.”
	}
}
{
%	\footnote{
%		your sheep and your oxen will stay put. Still negotiating, Pharaoh now
%		agrees to let the children go but wants the livestock to stay. It is not
%		as strong a guarantee to force Israel to return as holding the children
%		was, but it would mean that the Israelites would not have sufficient
%		food to make a long journey, or at least it would mean that the
%		Egyptians would get the livestock.
%	}
}

\Verse{25}
{
	\hRJE{וַיֹּאמֶר מֹשֶׁה גַּם אַתָּה תִּתֵּן בְּיָדֵנוּ זְבָחִים וְעֹלֹת וְעָשִׂינוּ לַיהֹוָה אֱלֹהֵינוּ׃}
}
{
	\eRJE{
		And Moses said, “You shall put sacrifices and offerings in
		our hand as well, so we’ll do them for YHWH, our God.
	}
}
{
}

\Verse{26}
{
	\hRJE{וְגַם מִקְנֵנוּ יֵלֵךְ עִמָּנוּ לֹא תִשָּׁאֵר פַּרְסָה כִּי מִמֶּנּוּ נִקַּח לַעֲבֹד אֶת יְהֹוָה אֱלֹהֵינוּ וַאֲנַחְנוּ לֹא נֵדַע מַה נַּעֲבֹד אֶת יְהֹוָה עַד בֹּאֵנוּ שָׁמָּה׃}
}
{
	\eRJE{
		And our livestock will go with us as well. Not a hoof will
		be left. Because we’ll take from it to serve YHWH, our
		God, and we won’t know how we’ll serve YHWH until we come
		there.”
	}
}
{
}

\Verse{27}
{
	\hRJE{וַיְחַזֵּק יְהֹוָה אֶת לֵב פַּרְעֹה וְלֹא אָבָה לְשַׁלְּחָם׃}
}
{
	\eRJE{
		And YHWH strengthened Pharaoh’s heart, and he was not
		willing to let them go.
	}
}
{
}

\Verse{28}
{
	\hRJE{וַיֹּאמֶר לוֹ פַרְעֹה לֵךְ מֵעָלָי הִשָּׁמֶר לְךָ אַל תֹּסֶף רְאוֹת פָּנַי כִּי בְּיוֹם רְאֹתְךָ פָנַי תָּמוּת׃}
}
{
	\eRJE{
		And Pharaoh said to him, “Go away from me! Watch yourself!
		Don’t continue to see my face. Because in the day you see
		my face you’ll die!” 10:28,29; 11:6. continue. The
		word—Hebrew root ysp, meaning to do again, add,
		continue—occurs three times in succession, a final
		reminder of Joseph (yôsp), whom the oppressing Pharaohs
		did not know.
	}
}
{
}

\Verse{29}
{
	\hRJE{וַיֹּאמֶר מֹשֶׁה כֵּן דִּבַּרְתָּ לֹא אֹסִף עוֹד רְאוֹת פָּנֶיךָ׃}
}
{
	\eRJE{
		And Moses said, “So you’ve spoken: I won’t continue to see
		your face anymore.”
	}
}
{
}
